\section{A commitment carries its credit history into the loan}
\label{sec:bank-commitment-draw}

A promise to lend can create a credit-loss provision before it creates a
loan or consumes cash. Once the customer draws, the uncertainty about whether
funding will occur has been resolved, but the customer's earlier credit history
has not disappeared. The distinction connects the commitment populations
reported above to the funded-loan method that follows. AASB 9.5.5.6 dates the
initial impairment comparison from the irrevocable commitment; B5.5.47 treats
the resulting financial asset as a continuation of that commitment
\citep{aasb9dated}.

Consider a synthetic bank with cash 200, deposits 180 and equity 20. It grants
an irrevocable commitment allowing one immediate draw of 100 or lapse. The
contract is fixed rather than revolving, cannot be net settled, is not
designated at fair value through profit or loss, and does not promise a
below-market loan. No fee or separate noncontingent commitment liability is
assumed. The bank will hold the resulting loan to collect principal and
interest. Its admitted initial fair value is 100, equal to its cash price,
and its payment of 121 after two years gives an effective interest rate of
10\%. These are explicit measurement and classification facts. The credit
probabilities below do not themselves establish market value; a different
initial value or a below-market commitment would require the separate
analysis under AASB 9.2.3, 4.2.1(d) and 5.1.1--5.1.1A.

All commitment and draw snapshots are ordered on 1 January 2024, before any
material interest accrues. Initially, the probability of a full draw is one
half. Conditional on drawing, default occurs in the first year with probability
0.1, in the second with probability 0.1, or does not occur with probability
0.8. These are mutually exclusive outcomes, not cumulative default
probabilities. Either default produces collection of 108.9 at the contractual
two-year maturity; no default produces 121. The present value of the cash
shortfall from either default is
\[
  \frac{121-108.9}{1.1^2}=10.
\]
The first-year default's loss is paid out at a later date, yet belongs in
twelve-month expected credit loss (ECL). The horizon selects possible default
events; it does not truncate their later cash shortfalls. Including expected
draws as required by B5.5.30--31 gives an initial provision of
\(0.5\times0.1\times10=0.5\). Debit impairment expense 0.5 and credit
the commitment provision 0.5. The nominal limit 100 is neither a funded asset
nor an impairment expense.

Before the draw, a new forecast and a separately supported significant
increase in credit risk (SICR) change the conditional default probabilities
to 0.2 in each year and 0.6 for no default. The borrower is not credit-impaired.
Doubling a probability is not a universal SICR rule; the credit finding is
a distinct input in this illustration. With draw probability still one half,
lifetime ECL is now \(0.5(0.2+0.2)10=2\). Debit a further impairment
expense of 1.5 and credit the provision by 1.5.

The customer now actually draws the whole 100. To see what changes, the
joint probabilities immediately before funding are 0.5 for no draw, 0.1
for draw and first-year default, 0.1 for draw and second-year default, and
0.3 for draw without default. Observing the draw removes the first outcome
and conditions the other three on their total probability of 0.5. The
funded loan therefore has lifetime ECL
\[
  \left(\frac{0.1}{0.5}+\frac{0.1}{0.5}\right)10=4.
\]
Keeping the old draw probability would still report 2, which measures the
earlier uncertain commitment. Resetting the funded loan to twelve-month ECL
would also give 2, now by omitting the second default window. Equal wrong
amounts can conceal different errors. The original commitment date, credit
baseline, current finding and actual draw must all survive the transition.

For a direct payment outside the bank, debit gross loan 100 and credit cash
100. In the same draw event, debit the existing provision 2, debit additional
impairment expense 2 and credit the loan allowance 4. The first 2 moves an
existing loss estimate from a liability into a deduction from the funded
asset; the other 2 reflects the information that funding has occurred.
AASB 7.B8E distinguishes commitment provisions from financial-asset allowances,
while 35H--35I and B8D require an explanation of their movements
\citep{aasb7dated}. Funding and its impairment transition are committed
together under the existing event procedure, even though the entries are
explained separately here.

There is no remaining undrawn capacity. Cash 100 plus gross loan 100 less
allowance 4 gives net assets 196, equal to deposits 180 plus equity 16.
Cumulative impairment expense is \(0.5+1.5+2=4\). Keeping the old provision
of 2 as well as the full funded allowance would count the loss twice and
leave equity 14. Using allowance 2 would instead leave equity 18. All three
sets of accounts can balance; only one measures the admitted funded exposure.

The route through the borrower's bank account has a different intermediate
state. Debit gross loan 100 and credit the borrower's deposit 100, with the
same provision-to-allowance transition. Cash stays 200, deposits become 280,
and equity is 16 after allowance 4. If the borrower then transfers 100 to
another bank, debit deposits and credit settlement cash by 100. The final
state matches direct disbursement. Creating the deposit did not itself
consume settlement funds; the later transfer did.

Alternatively, the customer lets the binding window expire without drawing.
Debit the provision 2 and credit impairment expense 2. Cash 200, deposits
180 and equity 20 remain, with no loan. This is an admitted lapse, not an
invented cancellation right. A partial or revolving draw needs separate
funded and remaining-exposure measurements. B8E allows combined allowance
presentation when their losses cannot be separately identified; any excess
over the gross funded asset is a provision. Neither that condition nor the
longer revolving exposure period in AASB 9.5.5.20 follows merely from calling
something a credit facility.

The probability tree is an authored assumption, not a calibrated forecast or
an observed default sample. Unknown probabilities can support conditional
losses but cannot produce a probability-weighted ECL without a measure.
Missing original credit history leaves staging unresolved, and missing
initial valuation leaves the funded measurement unresolved. The retained
AASB 9 compilation 5 and AASB 7 compilation 7 cover the stated historical
2024 case. Accounting equity of 16 does not determine regulatory capital:
the applicable jurisdiction, perimeter, deductions and exposure rules must
be supplied independently. The next example follows a funded loan through
interest, deterioration, collection and write-off using different stated facts.
